\documentclass[11pt]{article}
\usepackage[a4paper, top=18mm, bottom=18mm, left=20mm, right=20mm]{geometry}
\usepackage[T2A]{fontenc}
\usepackage[utf8]{inputenc}
\usepackage[ukrainian]{babel}
\usepackage{graphicx}
\setlength{\parindent}{0pt}
\setlength{\parskip}{4pt}
\pagestyle{empty}
\begin{document}
%begin
{\large\textbf{Контрольна робота}}

\textbf{Варіант \No\,1}\hfill \textit{ПІБ:}\,\underline{\hspace{60mm}}
\vspace{4mm}

%beginitem
1. Відмітити все, що є однією правильною лексемою:
( 1 ) 8,3	
( 2 ) <>	
( 3 ) 8.3	
( 4 ) FIFO

%enditem
%beginitem
2. Відмітити все, що є коректним оператором С++:
( 1 ) cout>>3;	
( 2 ) int x+y=z;	
( 3 ) cout<<3.4;	
( 4 ) int f(int n);

%enditem
%beginitem
3. Відмітити все, що є коректним оператором С++:
( 1 ) cout>>3;	
( 2 ) int x+y=z;	
( 3 ) cout<<3.4;	
( 4 ) int f(int n);

%enditem
%beginitem
4. Відмітити все, що є коректним оператором С++:
( 1 ) cout>>3;	
( 2 ) int x+y=z;	
( 3 ) cout<<3.4;	
( 4 ) int f(int n);

%enditem
%beginitem
5. Відмітити все, що є однією правильною лексемою:
( 1 ) 8,3	
( 2 ) <>	
( 3 ) 8.3	
( 4 ) FIFO

%enditem
%beginitem
6. Відмітити все, що є однією правильною лексемою:
( 1 ) 8,3	
( 2 ) <>	
( 3 ) 8.3	
( 4 ) FIFO

%enditem
%beginitem
7. Відмітити все, що є однією правильною лексемою:
( 1 ) 8,3	
( 2 ) <>	
( 3 ) 8.3	
( 4 ) FIFO

%enditem
%beginitem
8. Відмітити все, що є однією правильною лексемою:
( 1 ) 8,3	
( 2 ) <>	
( 3 ) 8.3	
( 4 ) FIFO

%enditem
%beginitem
9. Відмітити все, що є однією правильною лексемою:
( 1 ) 8,3	
( 2 ) <>	
( 3 ) 8.3	
( 4 ) FIFO

%enditem
%beginitem
10. Відмітити все, що є коректним оператором С++:
( 1 ) cout>>3;	
( 2 ) int x+y=z;	
( 3 ) cout<<3.4;	
( 4 ) int f(int n);

%enditem
%beginitem
11. Відмітити все, що є коректним оператором С++:
( 1 ) cout>>3;	
( 2 ) int x+y=z;	
( 3 ) cout<<3.4;	
( 4 ) int f(int n);

%enditem
%beginitem
12. Відмітити все, що є коректним оператором С++:
( 1 ) cout>>3;	
( 2 ) int x+y=z;	
( 3 ) cout<<3.4;	
( 4 ) int f(int n);

%enditem
%beginitem
13. Відмітити все, що є коректним оператором С++:
( 1 ) cout>>3;	
( 2 ) int x+y=z;	
( 3 ) cout<<3.4;	
( 4 ) int f(int n);

%enditem
%beginitem
14. Відмітити все, що є коректним оператором С++:
( 1 ) cout>>3;	
( 2 ) int x+y=z;	
( 3 ) cout<<3.4;	
( 4 ) int f(int n);

%enditem
%beginitem
15. Відмітити все, що є однією правильною лексемою:
( 1 ) 8,3	
( 2 ) <>	
( 3 ) 8.3	
( 4 ) FIFO

%enditem
%beginitem
16. Відмітити все, що є однією правильною лексемою:
( 1 ) 8,3	
( 2 ) <>	
( 3 ) 8.3	
( 4 ) FIFO

%enditem
%beginitem
17. Відмітити все, що є однією правильною лексемою:
( 1 ) 8,3	
( 2 ) <>	
( 3 ) 8.3	
( 4 ) FIFO

%enditem
%beginitem
18. Відмітити все, що є коректним оператором С++:
( 1 ) cout>>3;	
( 2 ) int x+y=z;	
( 3 ) cout<<3.4;	
( 4 ) int f(int n);

%enditem
%beginitem
19. Відмітити все, що є коректним оператором С++:
( 1 ) cout>>3;	
( 2 ) int x+y=z;	
( 3 ) cout<<3.4;	
( 4 ) int f(int n);

%enditem
%beginitem
20. Відмітити все, що є однією правильною лексемою:
( 1 ) 8,3	
( 2 ) <>	
( 3 ) 8.3	
( 4 ) FIFO

%enditem
%beginitem
21. Відмітити все, що є однією правильною лексемою:
( 1 ) 8,3	
( 2 ) <>	
( 3 ) 8.3	
( 4 ) FIFO

%enditem
%beginitem
22. Відмітити все, що є однією правильною лексемою:
( 1 ) 8,3	
( 2 ) <>	
( 3 ) 8.3	
( 4 ) FIFO

%enditem
%beginitem
23. Відмітити все, що є коректним оператором С++:
( 1 ) cout>>3;	
( 2 ) int x+y=z;	
( 3 ) cout<<3.4;	
( 4 ) int f(int n);

%enditem
%beginitem
24. Відмітити все, що є коректним оператором С++:
( 1 ) cout>>3;	
( 2 ) int x+y=z;	
( 3 ) cout<<3.4;	
( 4 ) int f(int n);

%enditem
%beginitem
25. Відмітити все, що є коректним оператором С++:
( 1 ) cout>>3;	
( 2 ) int x+y=z;	
( 3 ) cout<<3.4;	
( 4 ) int f(int n);

%enditem
%beginitem
26. Відмітити все, що є коректним оператором С++:
( 1 ) cout>>3;	
( 2 ) int x+y=z;	
( 3 ) cout<<3.4;	
( 4 ) int f(int n);

%enditem
%beginitem
27. Відмітити все, що є однією правильною лексемою:
( 1 ) 8,3	
( 2 ) <>	
( 3 ) 8.3	
( 4 ) FIFO

%enditem
%beginitem
28. Відмітити все, що є однією правильною лексемою:
( 1 ) 8,3	
( 2 ) <>	
( 3 ) 8.3	
( 4 ) FIFO

%enditem
%beginitem
29. Відмітити все, що є однією правильною лексемою:
( 1 ) 8,3	
( 2 ) <>	
( 3 ) 8.3	
( 4 ) FIFO

%enditem
%beginitem
30. Відмітити все, що є коректним оператором С++:
( 1 ) cout>>3;	
( 2 ) int x+y=z;	
( 3 ) cout<<3.4;	
( 4 ) int f(int n);

%enditem






























%end
\newpage
\end{document}

